%! Modeling with Exponential Functions: Interest, Half-Life, and e — Unit 6, Lesson 6.4 — Slides (Beamer)
\documentclass[aspectratio=169, 11pt]{beamer}
\usepackage{algebra2-beamer}   % provides \forestheader{} and \sectionlabel[color]{}
% Growth curves on money-sized grids. pgfmath's pow() is unreliable for fractional
% exponents, so every exponential is plotted as a*exp((ln b)t).
% ln(1.05) = 0.0487902   ln(0.5) = -0.693147

\begin{document}

% TITLE SLIDE (CourseName is not defined in beamer, so the name is literal)
{
\setbeamercolor{background canvas}{bg=forest}
\begin{frame}[plain]
  \vspace{2.8cm}\hspace{0.55cm}%
  \begin{minipage}{0.9\textwidth}
    {\color{goldacc}\bfseries\small LESSON 6.4}\\[0.5cm]
    {\color{white}\bfseries\LARGE Modeling: Compound Interest, Half-Life, and $e$}\\[1.2cm]
    {\color{forestmid}\small Algebra 2~$\cdot$~Unit 6: Exponential Functions}
  \end{minipage}
\end{frame}
}

% HOOK
\begin{frame}[t, plain]
  \forestheader{Four banks, one deposit, one rate --- and \$1.83}
  \sectionlabel{Hook}
  \begin{columns}[T]
    \begin{column}{0.52\textwidth}
      \footnotesize
      \renewcommand{\arraystretch}{1.35}
      \begin{tabular}{|l|c|c|}
      \hline
      \textbf{Bank} & \textbf{Interest added} & \textbf{After 1 year} \\
      \hline
      Ashmore   & annually  & \$1060.00 \\
      Bellhaven & quarterly & \$1061.36 \\
      Corliss   & monthly   & \$1061.68 \\
      Dunmore   & daily     & \$1061.83 \\
      \hline
      \end{tabular}
      \\[6pt]
      {\scriptsize \$1000 deposited, $6\%$ advertised at every bank.}
    \end{column}
    \begin{column}{0.44\textwidth}
      \footnotesize
      \begin{itemize}\setlength{\itemsep}{4pt}
        \item Dunmore compounds \textbf{365 times} as often as Ashmore --- and earns \$1.83 more.
        \item Annual $\to$ quarterly: $+\$1.36$. Monthly $\to$ daily: $+\$0.15$.
      \end{itemize}
      \vspace{3pt}
      \begin{block}{The question}
        Is there a bank aggressive enough to double your money in a year --- or is something stopping
        all of them?
      \end{block}
    \end{column}
  \end{columns}
\end{frame}

% THE SENTENCE
\begin{frame}[t, plain]
  \forestheader{One sentence runs the whole lesson}
  \sectionlabel[goldacc]{The rule}
  \begin{center}
    \Large The base is what happens in \textbf{one period}. \\[4pt]
    The exponent \textbf{counts the periods}.
  \end{center}
  \vspace{8pt}
  \begin{center}
    \Large $A = P\left(1+\dfrac{r}{n}\right)^{nt}$
  \end{center}
  \vspace{6pt}
  \begin{columns}[T]
    \begin{column}{0.48\textwidth}
      \begin{block}{Why $\frac{r}{n}$}
        \footnotesize One period earns only a fraction of the year's rate. At $8\%$ quarterly, a
        quarter pays $2\%$.
      \end{block}
    \end{column}
    \begin{column}{0.48\textwidth}
      \begin{block}{Why $nt$}
        \footnotesize The exponent is measured in \emph{periods}, not years. Five years of quarters is
        $20$ periods.
      \end{block}
    \end{column}
  \end{columns}
  \vspace{5pt}
  \begin{center}{\footnotesize This is still Lesson 6.1's $y=ab^{x}$ --- only the length of ``one
  period'' changed.}\end{center}
\end{frame}

% THE TWO ERRORS
\begin{frame}[t, plain]
  \forestheader{Break that sentence in half and you get today's two errors}
  \sectionlabel{Watch out}
  \begin{center}
    {\footnotesize \$1000 at $8\%$ compounded quarterly for $5$ years.}
  \end{center}
  \vspace{4pt}
  \begin{columns}[T]
    \begin{column}{0.32\textwidth}
      \begin{block}{Renata: $1000(1.08)^{20}$}
        \footnotesize Right count, unsplit rate. A full year's interest \emph{every quarter}.
        \\[3pt] Gives \$4660.96 --- the balance after $20$ \textbf{years}.
      \end{block}
    \end{column}
    \begin{column}{0.32\textwidth}
      \begin{block}{Miles: $1000(1.02)^{5}$}
        \footnotesize Right rate, unconverted time. Only five \emph{quarters} of growth.
        \\[3pt] Gives \$1104.08 --- the balance after $1\frac14$ years.
      \end{block}
    \end{column}
    \begin{column}{0.32\textwidth}
      \begin{block}{Correct: $1000(1.02)^{20}$}
        \footnotesize Base and exponent agree: a quarter's factor, counted in quarters.
        \\[3pt] $=\$1485.95$.
      \end{block}
    \end{column}
  \end{columns}
  \vspace{6pt}
  \begin{center}
    {\footnotesize Opposite halves of one mistake --- and they need different re-teaches.}
  \end{center}
\end{frame}

% e
\begin{frame}[t, plain]
  \forestheader{Chop the year finer and finer --- the numbers stop climbing}
  \sectionlabel[goldacc]{The number $e$}
  \begin{center}
    {\footnotesize \$1 at $100\%$ for one year, in $m$ equal pieces: $\left(1+\frac1m\right)^{m}$.}
    \\[6pt]
    \footnotesize
    \renewcommand{\arraystretch}{1.3}
    \begin{tabular}{|l|c|c|c|c|c|c|}
    \hline
    $m$ & $1$ & $10$ & $100$ & $1000$ & $10{,}000$ & $100{,}000$ \\
    \hline
    $\left(1+\frac1m\right)^{m}$ & $2$ & $2.59374$ & $2.70481$ & $2.71692$ & $2.71815$ & $2.71827$ \\
    \hline
    \end{tabular}
  \end{center}
  \vspace{8pt}
  \begin{columns}[T]
    \begin{column}{0.46\textwidth}
      \footnotesize The digits are freezing from the left. They are closing in on
      \[ e \approx 2.71828 \]
      an irrational number, like $\pi$ --- and like $\pi$, it has a button.
    \end{column}
    \begin{column}{0.50\textwidth}
      \begin{block}{Continuous compounding}
        \[ A = Pe^{\,rt} \]
        \footnotesize No $n$: the chopping has gone as far as chopping can go.
      \end{block}
    \end{column}
  \end{columns}
\end{frame}

% THE CEILING
\begin{frame}[t, plain]
  \forestheader{$e$ answers the Hook: the ceiling is low}
  \sectionlabel{Payoff}
  \begin{center}
    {\footnotesize \$2500 at $4\%$ for $6$ years --- the same deal, three schedules.}
  \end{center}
  \vspace{6pt}
  \begin{center}
    \footnotesize
    \renewcommand{\arraystretch}{1.4}
    \begin{tabular}{|l|c|c|}
    \hline
    \textbf{Compounded} & \textbf{Balance} & \textbf{Gained over the row above} \\
    \hline
    annually     & \$3163.30 & --- \\
    quarterly    & \$3174.34 & $+\$11.04$ \\
    continuously & \$3178.12 & $+\$3.78$ \\
    \hline
    \end{tabular}
  \end{center}
  \vspace{8pt}
  \begin{center}
    \begin{minipage}{0.84\textwidth}
      \begin{block}{So no bank can double your money by compounding faster}
        \footnotesize Continuous compounding is the most that rate can \emph{possibly} produce, and it
        beats quarterly by \$3.78. Compounding schedules are not where returns come from.
      \end{block}
    \end{minipage}
  \end{center}
\end{frame}

% HALF-LIFE
\begin{frame}[t, plain]
  \forestheader{Same sentence, different clock}
  \sectionlabel{Half-life \& doubling}
  \begin{center}
    \large
    $A(t)=A_{0}\left(\frac12\right)^{t/h}$ \qquad\qquad $A(t)=A_{0}\,2^{\,t/d}$
  \end{center}
  \vspace{6pt}
  \begin{columns}[T]
    \begin{column}{0.54\textwidth}
      \footnotesize
      \textbf{80 mg tracer, half-life 6 hours.}
      \renewcommand{\arraystretch}{1.25}
      \begin{tabular}{|l|c|c|c|c|c|}
      \hline
      $t$ (h) & $0$ & $6$ & $12$ & $18$ & $9$ \\ \hline
      $t/6$ & $0$ & $1$ & $2$ & $3$ & $1.5$ \\ \hline
      $A(t)$ & $80$ & $40$ & $20$ & $10$ & $28.28$ \\ \hline
      \end{tabular}
      \\[6pt]
      {\scriptsize The exponent is measured in \emph{halvings}. Dividing by $h$ is the unit
      conversion --- and $1.5$ halvings is a perfectly good number.}
    \end{column}
    \begin{column}{0.42\textwidth}
      \begin{block}{Two Lesson 6.0 facts, still true}
        \footnotesize Range $(0,80]$, horizontal asymptote $A=0$. \\[3pt]
        Half of something is never nothing --- the model never reaches zero.
      \end{block}
    \end{column}
  \end{columns}
\end{frame}

% 6.3 STILL WORKS
\begin{frame}[t, plain]
  \forestheader{Yesterday's method survives --- when the target cooperates}
  \sectionlabel{Solving}
  \begin{columns}[T]
    \begin{column}{0.47\textwidth}
      \footnotesize
      \textbf{When is $10$ mg left?}
      \begin{align*}
        80\left(\tfrac12\right)^{t/6} &= 10 \\
        \left(\tfrac12\right)^{t/6} &= \tfrac18 = \left(\tfrac12\right)^{3} \\
        \tfrac{t}{6} &= 3 \\
        t &= 18 \text{ hours}
      \end{align*}
      {\scriptsize $10$ is $80$ halved three times --- a common base exists.}
    \end{column}
    \begin{column}{0.47\textwidth}
      \footnotesize
      \textbf{When is $25$ mg left?}
      \begin{align*}
        \left(\tfrac12\right)^{t/6} &= \tfrac{5}{16} \\
        &\;\;?
      \end{align*}
      \vspace{-6pt}
      \begin{block}{Same drug, same half-life}
        \footnotesize $\frac{5}{16}$ is not a power of $\frac12$. The method reaches
        $\frac12,\frac14,\frac18,\frac1{16}$ of the start --- and nothing in between.
      \end{block}
    \end{column}
  \end{columns}
\end{frame}

% THE WALL
\begin{frame}[t, plain]
  \forestheader{When does \$1000 at $5\%$ double?}
  \sectionlabel[goldacc]{The wall}
  \begin{columns}[T]
    \begin{column}{0.40\textwidth}
      \centering
      \begin{tikzpicture}[x=0.20cm, y=0.0013cm]
        \draw[linegray!45, thin, xstep=2, ystep=500] (0,0) grid (20,3000);
        \draw[->, thick, charcoal] (0,0)--(21.5,0) node[above right, font=\tiny]{$t$};
        \draw[->, thick, charcoal] (0,0)--(0,3200) node[left, font=\tiny]{$A$};
        \foreach \t in {5,10,15,20} \node[below, font=\tiny] at (\t,0) {$\t$};
        \foreach \a in {1000,2000,3000} \node[left, font=\tiny] at (0,\a) {$\a$};
        \draw[very thick, forest, domain=0:20, samples=120, smooth]
          plot (\x,{1000*exp(0.0487902*(\x))});
        \draw[goldacc, very thick] (0,2000)--(20,2000);
        \fill[charcoal] (14,1979.9) circle (3pt);
        \fill[charcoal] (15,2078.9) circle (3pt);
      \end{tikzpicture}
    \end{column}
    \begin{column}{0.56\textwidth}
      \footnotesize
      \[ 1000(1.05)^{t}=2000 \;\longrightarrow\; (1.05)^{t}=2 \]
      \begin{itemize}\setlength{\itemsep}{3pt}
        \item $A(14)=\$1979.93$, $A(15)=\$2078.93$ --- the answer is in there.
        \item Step 1 of Lesson 6.3: rewrite both sides over one base. \textbf{$2$ is not a power of
              $1.05$.}
      \end{itemize}
      \vspace{3pt}
      \begin{block}{Not the same as no solution}
        \footnotesize $2^{x}=-3$ has \emph{no} answer --- a range fact. \\
        $(1.05)^{t}=2$ has an answer this method cannot name.
      \end{block}
    \end{column}
  \end{columns}
\end{frame}

% ROADMAP
\begin{frame}[t, plain]
  \forestheader{Where Unit 6 goes next}
  \sectionlabel{Connections}
  \textbf{Today:} the base is one period and the exponent counts periods --- which builds compound
  interest, continuous growth with $e$, half-life, and doubling time. Decay models keep Lesson 6.0's
  asymptote: they approach zero forever and never arrive.\\[8pt]
  \begin{itemize}\setlength{\itemsep}{3pt}
    \item \textbf{6.5} Exponential regression --- real data, no one to hand you the rate, and the
          fingerprint test from 6.0 decides the family.
    \item \textbf{Unit 7} Logarithms --- the tool that finally names every crossing point we could only
          bracket.
  \end{itemize}
  \vspace{6pt}
  \begin{center}
    \begin{minipage}{0.88\textwidth}
      {\footnotesize \textbf{The standing list:} $2^{x}=5$, \; $132(0.85)^{t}=32$, \;
      $(0.9)^{t}=2.118$, \; $(1.05)^{t}=2$, \; $(0.82)^{t}=0.5$. Every one has an answer. Not one has
      a common base.}
    \end{minipage}
  \end{center}
\end{frame}

\end{document}
